\section{Introduction}

Although camera-based human sensing systems are widely deployed, their practical
use can be limited by privacy concerns, illumination changes, and visual
occlusion. Wi-Fi sensing has therefore emerged as a promising non-contact and
non-visual alternative. By exploiting variations in wireless signal propagation
caused by human movement, Wi-Fi sensing can support human sensing applications
using existing communication infrastructure without requiring users to wear
dedicated sensors.

Most Wi-Fi sensing studies rely on Channel State Information (CSI). Compared with
coarse-grained measurements such as the Received Signal Strength Indicator
(RSSI), CSI provides subcarrier-level amplitude and phase responses across
transmit--receive antenna pairs \cite{halperin2011tool}. Variations in these
measurements have been used for a range of sensing applications, including
indoor localization \cite{kotaru2015spotfi}, respiration monitoring
\cite{wang2017tensorbeat}, gesture recognition \cite{zheng2019widar3}, and
human activity recognition \cite{wang2015understanding}.

Despite its sensing capability, CSI is not readily available from many commercial
Wi-Fi devices. Its extraction commonly depends on supported network interface
cards, modified firmware, monitor-mode configurations, or specialized CSI
collection tools \cite{gringoli2019free}. These requirements can restrict the
practical deployment of CSI-based sensing systems. Recent studies have therefore
investigated Beamforming Feedback Information (BFI) as an alternative source of
sensing information \cite{beamsense2023,bfmsense2024,beamsense2025,wibfi2023}.
In IEEE 802.11ac/ax networks, a station estimates the wireless channel from a
sounding transmission and returns compressed beamforming feedback to the access
point \cite{ieee2021ax,ieee2013ac}. Because BFI is produced as part of
standardized beamforming procedures, it provides an opportunity to perform
sensing using information already exchanged during Wi-Fi communication.

BFI represents wireless channel information in a compressed form. To reduce
feedback overhead, the beamforming matrix is encoded using quantized angular
parameters, commonly denoted by $\phi$ and $\psi$
\cite{itahara2022beamforming,li2024efficient}. These Beamforming Feedback Angles
(BFAs) describe the complex-valued beamforming matrix $\mathbf{V}$, which
corresponds to the right singular matrix of the estimated wireless channel and
represents the transmit-side spatial directions used for beamforming. Although
this compressed representation retains channel-dependent information affected by
human activity, its data structure and physical meaning differ from those of
conventional CSI. BFI should therefore be handled according to the characteristics
of its beamforming representation.

As with other data-driven wireless sensing approaches, the performance of
BFI-based activity recognition depends on the availability of sufficient labeled
training data. Collecting and labeling wireless measurements is time-consuming,
particularly when the measurements must cover multiple activities, participants,
and deployment environments. Conventional augmentation operations such as noise
injection, scaling, or element-wise mixing may increase the number of samples but
do not explicitly model the circular quantization and temporal evolution of BFAs.
This motivates a generative approach tailored to the structure used by downstream
BFI sensing models.

RF-Diffusion has demonstrated the potential of generative augmentation for
time-series RF signals by learning temporal, frequency-domain, and complex-valued
signal characteristics \cite{chi2024rfdiffusion}. Motivated by the accessibility
of BFI, we first explored extending this approach to BFI by using the reconstructed
complex beamforming matrix $\mathbf{V}$ as the generation target. However, being
complex-valued does not make $\mathbf{V}$ structurally equivalent to CSI. CSI
directly describes a channel response, whereas $\mathbf{V}$ represents channel
directions and is governed by normalization and geometric constraints as well as
phase ambiguity. Our exploratory results showed that direct complex-matrix
generation could preserve the expected tensor shape and basic numerical validity
while failing to reliably retain the temporal beam dynamics and activity semantics
required for BFI-based sensing. This mismatch indicates that extending a
CSI-oriented RF generator through input-format modification alone is insufficient
for BFI sensing augmentation.

To address this problem, we present \textbf{BeamDiff}, a sensing-aware diffusion
framework specialized for BFA sequence generation. Rather than independently
generating absolute BFA frames, BeamDiff preserves the first real frame as an
anchor and generates the nine frame-to-frame transitions of a ten-frame window.
It represents each transition as a shortest-path circular delta according to the
quantization range of each BFA channel, thereby avoiding artificial discontinuities
at angular boundaries. The denoising network is conditioned on the first-frame
anchor, diffusion timestep, and requested activity. In addition to the standard
diffusion noise-prediction objective, BeamDiff uses clean-delta reconstruction and
temporal-fidelity losses together with supervision from a frozen BeamSense
classifier. The generated transitions are accumulated from the anchor and mapped
back to valid quantized BFA values, producing synthetic samples that can be
directly consumed by existing BFI sensing models.

We evaluate BeamDiff using BFI measurements from the HAR-1 Kitchen and HAR-3
Classroom scenarios provided by the CSI-BFI-HAR dataset. We compare paired
real-only and real-plus-synthetic training protocols while keeping the real
validation and test folds fixed and augmenting only the training folds. This
evaluation examines generated-data fidelity, downstream sensing performance,
and robustness across model initialization seeds.

The main contributions of this work are summarized as follows:

\begin{itemize}
    \item We investigate the extension of complex RF generative augmentation to
    BFI and identify a structural mismatch between generic complex-matrix
    generation and the beam-direction dynamics required for BFI-based sensing.

    \item We propose BeamDiff, a BFA-specific conditional diffusion framework
    that generates shortest-path circular angle transitions from a real
    first-frame anchor rather than independently generating quantized frames.

    \item We introduce a sensing-aware objective that combines diffusion noise
    prediction, clean-delta reconstruction, temporal-fidelity regularization, and
    frozen-classifier supervision to preserve both BFA dynamics and activity
    semantics.

    \item We conduct paired downstream evaluations in two BFI sensing
    environments and analyze augmentation gain, generated-data fidelity, and
    seed-dependent behavior.
\end{itemize}

% Before final submission, verify all BibTeX keys and report detailed numerical
% results only in the abstract and evaluation section.
